\chapter{Budgeted Lifting Through Oracles}
\label{ch:budgeted}

\section{Adversarial Query Budgets}

ROM proofs depend on query counts.  The development therefore proves that hash
and signing query budgets are preserved by the relevant wrappers.  The budgeted
interfaces certify that prequery and programming losses are multiplied by the
right quantities.

\section{Pen-and-Paper Fundamental-Lemma Lifting}

\begin{lemma}[Budgeted one-call lifting]
Suppose two signing-oracle implementations \(O_0\) and \(O_1\) are identical
unless a clean event \(C\) fails during one signing call.  If for every state
satisfying the budget invariant
\[
\Pr[\neg C\text{ during one call}]\le \epsilon,
\]
and the adversary makes at most \(q_S\) signing calls, then the distinguishing
loss between the two full games is at most \(q_S\epsilon\), plus any separately
accounted hash-query bad events.
\end{lemma}

\begin{proof}
Index the signing calls from \(1\) to \(q_S\).  Define a hybrid sequence in
which the first \(i\) calls use \(O_1\) and the remaining calls use \(O_0\).  By
the fundamental lemma, adjacent hybrids differ by at most the probability that
the clean event fails in call \(i+1\).  The budget invariant ensures that each
such probability is bounded by \(\epsilon\).  Summing over at most \(q_S\) calls
gives \(q_S\epsilon\).  Hash-query bad events are added by another union-bound
or fundamental-lemma step.
\end{proof}

\section{One-Call Signing Loss}

The concrete one-call signing loss relates an attempt-side clean projection to
an exact-side projection.  Lemmas such as
\code{concrete\_budgeted\_ro\_signing\_attempt\_o\_sign\_active\_clean\_signature\_structural\_le},
\code{concrete\_budgeted\_ro\_exact\_hyperball\_o\_sign\_active\_signature\_exact},
and
\code{concrete\_budgeted\_o\_sign\_clean\_one\_call\_loss\_from\_self\_logged\_surface}
encode this local proof.

\section{Detailed EasyCrypt Code Analysis}

The module \code{ROMInternalTranscriptBudgetedPaperSimAsNMA} is the central
budgeted wrapper.  Its signing oracle \code{O.sign} is where a one-step sampler
loss must be lifted through adversarial interaction.  The code first proves
projection lemmas: the attempt side projects from \code{O.sign} to the
self-logged \code{sample\_with\_seed} event, and the exact side projects from
\code{sample\_with\_seed} back to exact signing.

The lemma
\code{concrete\_budgeted\_self\_logged\_signature\_clean\_event\_loss\_surface}
acts as a reusable surface theorem.  It abstracts over the local event and
requires both projections.  The final lemma
\code{concrete\_budgeted\_o\_sign\_clean\_one\_call\_loss\_from\_self\_logged\_surface}
instantiates that surface for the concrete \haetae{} lazy ROM.  This structure
is stronger than a coarse local wrapper bound because it proves that the exact
events used by the signing oracle are the events bounded by the sampler proof.

\section{From One Call to NMA Wrapper Obligations}

The one-call bound supplies the local surface that a non-vacuous NMA-level lift
should consume.  The checked lemmas whose names contain
\code{rom\_internal\_budgeted\_nma} expose the bridge from local signing facts to
adversarial-game facts, including the concrete clean event and the
\code{sampler\_bad\_prequery} log.  The present top-level path still has coarse
fallback steps: some NMA bridge lemmas use probability bounds whose usefulness
depends on the oversized \code{rejection\_sampling\_loss\_term}.  The remaining
proof obligation is to make the final NMA composition consume the local clean
one-call surface directly.
